\section{Prime Skeleton: From Cusp Evaluation to Hecke Rigidity Verification}

\subsection{$q$-Expansion at Cusps and Finite-Truncation Auditability}

Mapping scanning data to modular parameter $\tau$ and performing $q=e^{2\pi i\tau}$ expansion in the cusp region yields coefficient data from standard objects like Eisenstein series, $\Delta(\tau)$, and $j(\tau)$. Since numerical sensitivity is amplified near cusps, any finite truncation must be accompanied by explicit tail bounds and stability estimates to guarantee ``auditable'' error control.

\subsection{Hecke Operators and Prime-Generated Recursive Constraints}

The Hecke algebra is generated by operators $T_p$ at primes; for eigenforms, Fourier coefficients satisfy strong constraints generated by prime recursion. At the finite data level, Sturm-type bounds can elevate ``finite-order agreement'' to full agreement: when coefficients satisfy Hecke verification within a certain range, deviations can be regarded as noise/protocol mismatch and their residuals quantified. This ``prime skeleton'' mechanism provides an arithmetic rigidity verification layer for end-to-end protocols, making scan---readout results not only statistically controllable but also verifiable in arithmetic structure.

\subsection{End-to-End Construction Chain}

Combining the above elements yields the following construction chain (all closed within operational definitions):
\[
\begin{aligned}
&\text{Scanning orbit}\ \to\ \text{Finite-resolution window readout}\ \to\ \text{Discrepancy certificate}\ \to\ \text{Cusp evaluation and $q$-coefficients}\\
&\to\ \text{Hecke/prime skeleton verification}\ \to\ \text{Zeckendorf/Ostrowski tick encoding and folding stable sectors}.
\end{aligned}
\]
This chain provides a unified protocol perspective spanning dynamical systems, numerical analysis, and arithmetic constraints: each step has explicit input/output and error bounds, and contributions can be isolated through control experiments.
